\documentclass[10pt,a4paper]{amsart}
\usepackage[margin=19mm]{geometry}
\usepackage{amsmath,amssymb,mathtools}
\usepackage[scr=rsfs,cal=euler]{mathalfa}
\usepackage{xfrac}
\usepackage[unicode,hidelinks,bookmarks=false]{hyperref}
\newcommand{\ie}{i.e.\ }
\newcommand{\cf}{cf.\ }
\newcommand{\resp}{respectively\ }
\newcommand{\Subsection}{\subsection}
\providecommand{\nobreakdash}{\nobreak\hbox{-}}
\setlength{\emergencystretch}{2em}
\multlinegap0pt
\usepackage{bm}
\usepackage{bbm}
\usepackage{dsfont}
\usepackage{xspace}
\usepackage{cancel}
\usepackage{mathtools}
\usepackage{amsthm}
\makeatletter
\@ifundefined{theorem}{\newtheorem{theorem}{Theorem}}{}
\@ifundefined{corollary}{\newtheorem{corollary}[theorem]{Corollary}}{}
\@ifundefined{definition}{\newtheorem{definition}[theorem]{Definition}}{}
\@ifundefined{lemma}{\newtheorem{lemma}[theorem]{Lemma}}{}
\makeatother
\title[Characterization of bi-parametric potentials and rate of con]{Characterization of bi-parametric potentials and rate of convergence of truncated hypersingular integrals in the Dunkl setting}
\author{Sandeep Kumar Verma, Athulya P}
\date{Source: arXiv:2505.15748v3 (CC BY 4.0)}
\begin{document}
\maketitle

\noindent\emph{Source and licence.} Characterization of bi-parametric potentials and rate of convergence of truncated hypersingular integrals in the Dunkl setting. Version arXiv:2505.15748v3 (CC BY 4.0), distributed under the Creative Commons Attribution 4.0 International licence, as recorded in the arXiv licence metadata for this exact version and shown on the abstract page https://arxiv.org/abs/2505.15748v3. Full rights notice: licenses/arxiv-2505.15748-CC-BY-4.0.txt.\par\smallskip

\noindent\textit{Editorial scope.} Counted unit: the Theorem whose source label is \texttt{Riez-R-of-C} in the article (source0.tex lines 931--944, source label \texttt{Riez-R-of-C}), delivered together with the article's own complete original proof of it (source0.tex lines 945--1016, one \texttt{proof} environment of 72 lines). No mathematics is written, repaired or re-derived: statement and proof are the article's own text line for line. Required context retained uncounted: beta-convolution (lines 383--385); Lemma-R-2 (lines 507--512); smoothness-1 (lines 819--825); corollary-5.4 (lines 858--869). Every definition, display and label of those lines is retained unchanged. Omitted from this package and not redistributed: 1--382, 386--506, 513--818, 826--857, 870--930, 1017--1161. Nothing inside a retained excerpt is omitted or altered: every retained body is delivered exactly as the article's lines are written, so no omission receipt of any kind accompanies this package. Citations: the retained excerpts carry no citation command at all, so no bibliography entry of the article is transcribed and no reference is redistributed. Reference adaptation: none was needed and none was made; every reference inside the delivered unit resolves inside it, so no unresolved reference or citation is produced. Printed slips kept literal: no printed word, symbol, hypothesis or number of the counted statement or of its proof was altered. Layout and class: the article's own class is \texttt{amsproc} with packages including amsmath, amsthm, amscd, amsfonts, amssymb, graphicx, color, subcaption, enumerate, hyperref, natbib; the wrapper is the lane's pinned \texttt{amsart} editorial wrapper at 10pt on A4, which restates the theorem-like environments the retained text needs and adds the article's own macro definitions that the retained text needs (none), with extra wrapper packages (bm, bbm, dsfont, xspace, cancel, mathtools). The delivered numbering is the unit's own. Assets: no retained span contains a figure environment, a graphics-inclusion command, a PDF-page inclusion, a table or a TikZ picture; no image file is copied into this package, the package declares no asset dependency and no image file is redistributed with it. Licence: version arXiv:2505.15748v3 of this article is distributed under the Creative Commons Attribution 4.0 International licence, as recorded in the arXiv licence metadata for this exact version and shown on the abstract page https://arxiv.org/abs/2505.15748v3. A full rights notice is delivered at licenses/arxiv-2505.15748-CC-BY-4.0.txt. Characters: every retained excerpt is pure ASCII, so the delivered engine, XeLaTeX with its default Latin Modern text font, reports no missing character.
\begin{align} \label{beta-convolution}
    \mathfrak{B}_k^{(\beta,t)}(f)(x)= \mathcal{W}_k^{(\beta,t)} \underset{k}{\ast} f(x),
\end{align}

\begin{lemma}\label{Lemma-R-2}
    Let $0< \alpha < n+2\gamma$, $\beta >\alpha $, $f\in L_k^p(\mathbb{R}^n)$ with $1\leq p < \frac{n+2\gamma}{\alpha}$ and $g=\mathcal{I}_k^\alpha(f)$. Then
    \begin{align*}
        \mathfrak{T}_\epsilon (g)(\xi)= \int_0^\infty \mathcal{B}_k^{(\beta, \epsilon s)}(f)(\xi)\mathcal{K}_{\alpha/\beta}(s)ds,
    \end{align*} where $\mathcal{K}_{\theta}(s)= \frac{1}{s} \frac{1}{\Gamma(\theta+1)}\int_0^s (s-r)^\theta d\nu(r)$, with $\theta=\alpha/\beta$. 
\end{lemma}

\begin{definition}
Let $0< \rho < 1$ and $B[0,r]$ be the closed ball around the origin with radius $r$.  We say that  $f\in L_k^{2}(\mathbb{R}^n)$ has $\eta$-smoothness property at $x_0$ if 
       \begin{align} \label{smoothness-1}
            \mathcal{N}_{\eta}f(x_0) = \underset{0<r\leq \rho}{\sup} \, \frac{1}{r^{n+2\gamma}\eta(r)} \int_{B[0,r]} |\tau_{x}f(x_0)-f(x_0)|w_k(x)dx < \infty.
       \end{align} 

\end{definition}

\begin{corollary} \label{corollary-5.4}
    Let $f\in L_k^2(\mathbb{R}^n)$ has $\eta$-smoothness property at $x_0$ and we denote $g(x)=|\tau_xf(x_0)-f(x_0)|$. For $0<t\leq 1, \beta \geq 1,$ and 
    $\phi_t^\beta(|x|)=\mathcal{W}_k^{(\beta,t)}(x)$, we have 
    \begin{align*}
       \left| \int_{B[0,\rho]} g(x)\phi_t^\beta(|x|) w_k(x)dx \right| \leq   
       C^1\eta(t^{\frac{1}{\beta}}) 
    \end{align*} and for  $0<\beta\leq 1$ we have 
    \begin{align*}
       \left| \int_{B[0,\rho]} g(x)\phi_t^\beta(|x|) w_k(x)dx \right| \leq C^{*}\eta(t^{\beta}), 
    \end{align*}
    where $C^1$ and $C^{*}$ are the constants independent of $t$. 
\end{corollary}

\begin{theorem}\label{Riez-R-of-C}
Let \( f \in L_k^p(\mathbb{R}^n) \cap L_k^2(\mathbb{R}^n) \) for \( 1 \leq p < \infty \), and suppose that \( f \) possesses \(\eta\)-smoothness at the point \( x_0 \in \mathbb{R}^n \). Then the following estimate holds:
\begin{align*}
    \left| \mathfrak{T}_\epsilon(\mathcal{I}_k^{\alpha}(f))(x_0)- f(x_0) \right| \leq C\,\eta(Y(\epsilon)) \quad \text{as } \epsilon \to 0^+,
\end{align*}
where \( 0 < \alpha < \frac{n+2\gamma}{p} \), \( C>0\) is a constant independent of \( \epsilon\), and the function \( Y(\epsilon) \) is defined by
\begin{equation*}
     Y(\epsilon) =   
\begin{cases} 
  \epsilon^{1/\beta}, & \text{if } \beta \geq 1, \\ 
  \epsilon^{\beta},   & \text{if } 0 < \beta \leq 1.
\end{cases}
\end{equation*}
\end{theorem}

\begin{proof}
We begin the proof by considering  $\mathfrak{T}_\epsilon(\mathcal{I}_k^{\alpha}(f))(x_0)- f(x_0)$. Invoking Lemma \ref{Lemma-R-2}, we have 
\begin{align} \label{m-eq0}
  i_1=  \left| \mathfrak{T}_\epsilon(\mathcal{I}_k^{\alpha}(f))(x_0)- f(x_0) \right| \leq \frac{1}{\left|C(\frac{\alpha}{\beta},\nu)\right|}\int_0^\infty \left|\mathfrak{B}_k^{(\beta, \epsilon s)}f(x_0)-f(x_0)\right| \left|\mathcal{K}_{\alpha/\beta}(s)\right|ds .
\end{align}
    Let us denote $i_2=| \mathfrak{B}_k^{(\beta, \epsilon s)}f(x_0)-f(x_0)|$. Using \eqref{beta-convolution}, we obtain
    \begin{align*}
        i_2 &= \left|\int_{\mathbb{R}^n} \mathcal{W}_k^{(\beta, \epsilon s)}(y) \left( \tau_yf(x_0)-f(x_0) \right)w_k(y)dy
        \right|\\
        &\leq \left|\int_{B[0,\rho]} \mathcal{W}_k^{(\beta, \epsilon s)} (y)\left( \tau_yf(x_0)-f(x_0) \right)w_k(y)dy
        \right| + \\ & \hspace{4.5cm}  \left|\int_{B[0,\rho]^c} \mathcal{W}_k^{(\beta, \epsilon s)}(y) \left( \tau_yf(x_0)-f(x_0) \right)w_k(y)dy
        \right|\\
        & = i_3 +i_4.
    \end{align*}
    In order to estimate $i_3$, we use Corollary \ref{corollary-5.4}. 
    \begin{align*}
      i_3=  \int_{B[0,\rho]} \mathcal{W}_k^{(\beta, \epsilon s)}(y) \left| \tau_yf(x_0)-f(x_0) \right|w_k(y)dy \leq C_1\eta((\epsilon s)^{1/\beta}) \leq C_1 (1+s^{1/\beta}) \eta (\epsilon^{1/\beta}),
    \end{align*} where $ 1\leq \beta \leq 2$,
    and $ i_3 \leq  C^\prime (1+s^{\beta}) \eta (\epsilon^{\beta}), \quad \text{ for $0< \beta \leq 1$. }$\\
\noindent  For $\beta > 2$, invoking \eqref{smoothness-1} and the bound of the function $ \mathcal{W}_k^{(\beta, \epsilon s)}(y), $ we have 
    \begin{align*}
        i_3 & \leq \int_{B[0,\rho]} \left|\mathcal{W}_k^{(\beta, \epsilon s)}(y) \left( \tau_yf(x_0)-f(x_0) \right) \right|w_k(y)dy 
         \leq C_2\epsilon s \mathcal{N}_{\eta}f(x_0) \rho^{n+2\gamma}\eta(\rho)\\ &\leq C_3s\eta(\epsilon^{1/\beta}), \quad \text{where} \quad C_3= \frac{1}{a} \mathcal{N}_\eta f(x_0)\rho^{n+2\gamma}\eta(\rho)C_2.
    \end{align*}

To calculate $i_4$, we decompose the integrals into two parts as
\begin{align*}
    i_4 &\leq \left|\int_{B[0,\rho]^c} \mathcal{W}_k^{(\beta, \epsilon s)}(y)  \tau_yf(x_0) w_k(y)dy
        \right| + \left|\int_{B[0,\rho]^c} \mathcal{W}_k^{(\beta, \epsilon s)}(y)  f(x_0) w_k(y)dy 
        \right| \\
        &= i_5+i_6.
\end{align*}
In view of the properties of convolution and translation, $i_5$ becomes
\begin{align*}
    i_5 &= \left|\int_{B[0,\rho]^c} \mathcal{W}_k^{(\beta, \epsilon s)}(y)  \tau_yf(x_0) w_k(y)dy
        \right| = \left| \int_{\mathbb{R}^n} \chi_{B[0,\rho]^c} \mathcal{W}_k^{(\beta, \epsilon s)}(y)\tau_{x_0}( f)(y)w_k(y)dy \right| \notag \\
        & \leq  \int_{\mathbb{R}^n} \left| \mathcal{W}_k^{(\beta, \epsilon s)}\right|(y)|\tau_{x_0}(f)(y)|w_k(y)dy. \notag
\end{align*} 
Using the Cauchy-Schwarz inequality, we obtain  
\begin{align*}
    i_5 \leq  \|\mathcal{W}_k^{(\beta, \epsilon s)}\|_{L_k^2(\mathbb{R}^n)} \|f\|_{L_k^2(\mathbb{R}^n)} \leq \epsilon s C_4 \leq C_5 s\eta(\epsilon^{1/\beta}), \quad \text{ for }  \beta \geq 1. %\label{i_5-1}
\end{align*}
In addition, 
\begin{align*}%\label{i_5-2}
    i_5 \leq \epsilon s C_4 \leq C^{\prime\prime} s\eta(\epsilon^{\beta}), \quad \text{ for } 0< \beta \leq 1.
\end{align*}

\noindent Moreover 
\begin{align*}
    i_6 \leq C_6 s\eta(\epsilon^{1/\beta}) \quad \text{ for $\beta \geq 1$  and } i_6 \leq C^{\prime\prime\prime} s\eta(\epsilon^{\beta}) \quad \text{ for $0<\beta \leq 1$.}
\end{align*}

\noindent We finally define 
 $C^1= \max\{C_1, C_3, C_6 \}$, and $C^*=\max \{C^\prime, C^{\prime\prime},C^{\prime\prime\prime} \}$. Then
\begin{align} \label{m-eq}
    i_2 \leq C^1\eta(\epsilon^{1/\beta})\left( 1+2s+s^{1/\beta} \right) \quad \text{for } \beta\geq 1,
\end{align} and 
\begin{align} \label{m-eq-2}
    i_2 \leq C^*\eta(\epsilon^{\beta})\left( 1+2s+s^{\beta} \right) \quad \text{for }0< \beta\leq 1.
\end{align}
 Invoking \eqref{m-eq} and \eqref{m-eq-2} in \eqref{m-eq0},  we have
\begin{align*}
     i_1  \leq C^1\eta(\epsilon^{1/\beta}) \int_0^\infty  \left( 1+2s+s^{1/\beta} \right)\left|\mathcal{K}_{\alpha/\beta}(s)\right|ds \leq C_1(\nu)\eta(\epsilon^{1/\beta}) \text{  for } \beta\geq 1,
\end{align*} 
and 
\begin{align*}
     i_1 \leq C^*\eta(\epsilon^{\beta}) \int_0^\infty  \left( 1+2s+s^{\beta} \right)\left|\mathcal{K}_{\alpha/\beta}(s)\right|ds \leq C_2(\nu)\eta(\epsilon^{\beta}) \text{ 
 for } 0<\beta\leq 1,
\end{align*}
where the properties of the wavelet measure $\nu$ guarantee the convergence of the integrals. Thus, we have the desired result by choosing constant \( C =\max\{C_1(\nu),C_2(\nu)\}\).
% $\left| \mathfrak{T}_\epsilon(\mathcal{I}_k^{\alpha}f)(x_0)- f(x_0) \right| \leq C(\nu)\eta(\epsilon^{1/\beta}),$ as $\epsilon \rightarrow 0^+.$
\end{proof}

\end{document}
